% ============================================
% MEM-01
% Gerado de molambudos.md — build_miolo.py
% ============================================
\fragtitle{mem}{M}{MEM}{01}{O Sertão Antes do Fim}
\markboth{MEM-01 --- O Sertão Antes do Fim}{}

\noindent \textbf{1914 --- Quixeramobim, Ceará}

\vspace{0.5em}\noindent\rule{\textwidth}{0.5pt}\vspace{0.5em}

\noindent Meu pai dizia que a terra tinha um jeito de avisar.

\noindent Primeiro o gado morre. Eles sentem antes da gente. O gado começa a cair --- não de doença, de fraqueza. O pasto some, o gado emagrece, os ossos aparecem por baixo do couro, e um dia o boi deita e não levanta mais. Meu pai matava os que ainda estavam em pé para não perder a carne. Mas não tinha água pra salgar. A carne apodrecia antes de a gente comer.

\noindent O calor não dava trégua nem de noite. A poeira entrava pela porta e cobria tudo com uma camada fina, cor de osso. Eu sentia o cheiro de carne estragada chegar antes do vento --- um cheiro doce e podre que, anos depois, eu reconheceria de novo, vindo de dentro de mim mesmo.

\noindent Depois do gado, o rio seca. O Quixeramobim era um rio que cantava. Eu ouvia ele da porta de casa --- um barulho de água nas pedras, baixinho, como quem conversa sozinho. Em 1914, ele parou de cantar. Primeiro ficou raso. Depois ficou barrento. Depois ficou seco. O leito parecia uma estrada de osso branco.

\noindent Depois do rio, as pessoas mudam.

\noindent Meu pai mudou primeiro. Ele começou a ficar horas sentado na soleira da porta, olhando pro nada. Minha mãe perguntava o que ele tinha. Ele dizia: \textit{\textquotedblleft{}Tô esperando.\textquotedblright{}} Não dizia o quê. Acho que ele mesmo não sabia.

\noindent Depois mudaram os vizinhos. Seu Lázaro, que era o homem mais alegre da região, parou de falar. Dona Sinhá, que fazia pão pra todo mundo, começou a guardar comida debaixo da cama. As crianças pararam de brincar.

\noindent Quando a comida acabou de vez, ninguém se despediu.

\noindent A gente simplesmente saiu.

\noindent Minha mãe arrumou um saco com três peças de roupa, uma cuia, e uma imagem de Nossa Senhora das Dores que era da avó dela. Meu pai não arrumou nada. Ele olhou a casa uma última vez. A porta aberta. A mesa vazia.

\noindent \textit{\textquotedblleft{}Joaquim\textquotedblright{}\textit{, ele disse. }\textquotedblleft{}Vai olhar pra frente. Não olha pra trás.\textquotedblright{}}

\noindent Eu olhei pra trás.

\noindent A casa estava queimando.

\noindent O fogo era a luz do sol poente que pegava no telhado de palha e fazia a casa parecer em chamas. Parecia que o sertão inteiro estava pegando fogo.

\noindent Eu pensei: \textit{a gente não vai voltar.}

\noindent A gente não voltou.

\vspace{0.5em}
\noindent\textit{\textbf{↪ Rotas:} \rota{MEM-02} \rotasep \rota{DOC-02}}